\section{Results Discussion}
\label{sec:results}

\subsection{QR Factorization and Linear Least-Squares Solution}
\label{subsec:problem1-results}

The modified Gram--Schmidt algorithm was applied to the
overdetermined system defined in Section~\ref{sec:problem}. The
computed factors satisfy $A\approx QR$, where
$Q\in\mathbb{R}^{5\times4}$ has approximately orthonormal columns
and $R\in\mathbb{R}^{4\times4}$ is upper triangular.

The computed matrices, rounded to six decimal places, are

\begin{equation}
    Q\approx
    \begin{bmatrix}
        0.295599  & 0.242348  & -0.209363 & 0.159387  \\
        0.197066  & -0.200200 & 0.322900  & 0.898629  \\
        0.492665  & 0.140811  & 0.781674  & -0.352688 \\
        -0.098533 & 0.938737  & -0.000348 & 0.202929  \\
        -0.788263 & 0.011495  & 0.490803  & 0.038631
    \end{bmatrix},
    \label{eq:computed-Q}
\end{equation}

and

\begin{equation}
    R\approx
    \begin{bmatrix}
        10.148892 & -4.926646 & -3.251587 & -4.236916 \\
        0         & 10.135490 & 2.859317  & 4.846950  \\
        0         & 0         & 13.048045 & 2.480397  \\
        0         & 0         & 0         & 10.020142
    \end{bmatrix}.
    \label{eq:computed-R}
\end{equation}

The quality of the factorization was assessed using the relative
reconstruction error and the orthogonality error. The computed
values were

\begin{equation}
    \eta_{\mathrm{QR}}=1.0683\times10^{-16},
    \qquad
    \eta_Q=4.5180\times10^{-16}.
\end{equation}

Both indicators are close to the expected scale of floating-point
rounding errors for this computation. These results support the
numerical accuracy of the factorization and the orthogonality of
the computed columns of $Q$.

The least-squares solution was obtained by solving
$R\mathbf{x}^{*}=Q^T\mathbf{b}$ using back substitution. The
resulting solution, rounded to nine decimal places, is

\begin{equation}
    \mathbf{x}^{*}\approx
    \begin{bmatrix}
        -0.404938762 \\
        -0.040220457 \\
        0.992366790  \\
        -0.568506366
    \end{bmatrix}.
    \label{eq:computed-solution}
\end{equation}

The Euclidean norm of the residual was

\begin{equation}
    \|\mathbf{A}\mathbf{x}^{*}-\mathbf{b}\|_2
    \approx 7.677354.
\end{equation}

The nonzero residual indicates that the original system is
inconsistent. This is expected for an overdetermined least-squares
problem and does not, by itself, indicate a failure of the numerical
method. Instead, the solution minimizes the residual norm among
all vectors in $\mathbb{R}^4$.

To assess the optimality of the computed solution, the norm of
the residual projected onto the column directions of $A$ was also
evaluated. The result was

\begin{equation}
    \left\|A^T(A\mathbf{x}^{*}-\mathbf{b})\right\|_2
    \approx 6.1020\times10^{-14}.
\end{equation}

This value is numerically small and is consistent with the
least-squares optimality condition $A^T(A\mathbf{x}^{*}-\mathbf{b})
    =\mathbf{0}$. Hence, the computed residual is approximately
orthogonal to the column space of $A$.

Finally, the diagonal entries of $R$ were compared with the
scale-dependent numerical tolerance used by the implementation.
All four diagonal entries were substantially larger than this
tolerance, supporting the numerical assessment that $A$ has full
column rank. Therefore, the least-squares solution is unique.

These results validate the use of the modified Gram--Schmidt
factorization and back substitution for the matrix considered in
this experiment. Nevertheless, modified Gram--Schmidt without
column pivoting is not a universally reliable rank-revealing
method for severely ill-conditioned matrices. Householder QR
or rank-revealing factorizations may be preferable when stronger
numerical robustness is required
\cite{golub2013matrix,higham2002accuracy}.

\subsection{Quadratic Least-Squares Fitting}
\label{subsec:quadratic-fitting}

The second experiment considered the least-squares approximation
of six observations by a polynomial of degree at most two.
The design matrix $X\in\mathbb{R}^{6\times3}$ was constructed
using the columns associated with $x^2$, $x$, and the constant
term. The coefficient vector
$\boldsymbol{\theta}=[a,b,c]^T$ was computed using the reduced
QR factorization and back substitution.

The computed coefficients were
\begin{equation}
    \boldsymbol{\theta}^{*}\approx
    \begin{bmatrix}
        1.9739\times10^{-15} \\
        0.4000               \\
        0.6000
    \end{bmatrix}.
    \label{eq:quadratic-coefficients}
\end{equation}

Thus, the fitted polynomial is
\begin{equation}
    p(x)\approx 0.4000x+0.6000,
    \label{eq:fitted-polynomial}
\end{equation}
because the estimated quadratic coefficient is numerically
negligible. Although a quadratic model was specified, the resulting
least-squares polynomial is effectively linear at the displayed
precision. This behavior is a property of the data and does not
indicate a failure of the numerical algorithm.

Table~\ref{tab:quadratic-fit-results} compares the observations
with the fitted values and the corresponding residuals.

\begin{table}[htbp]
    \centering
    \small
    \caption{Observations, fitted values, and residuals.}
    \label{tab:quadratic-fit-results}
    \begin{tabular}{rrrr}
        \hline
        $x_i$ & $y_i$ & $\widehat{y}_i$ & $y_i-\widehat{y}_i$ \\
        \hline
        1     & 1     & 1.0             & 0.0                 \\
        2     & 1     & 1.4             & -0.4                \\
        3     & 3     & 1.8             & 1.2                 \\
        4     & 1     & 2.2             & -1.2                \\
        5     & 3     & 2.6             & 0.4                 \\
        6     & 3     & 3.0             & 0.0                 \\
        \hline
    \end{tabular}
\end{table}

The residual vector has Euclidean norm
\begin{equation}
    \|\mathbf{y}-X\boldsymbol{\theta}^{*}\|_2
    \approx 1.7889.
\end{equation}
The mean squared error is
\begin{equation}
    \operatorname{MSE}
    =\frac{1}{6}\sum_{i=1}^{6}
    (y_i-\widehat{y}_i)^2
    =\frac{3.2}{6}
    \approx 0.5333.
    \label{eq:quadratic-mse}
\end{equation}

The residual is nonzero because the fitted polynomial does not
pass through every observation. This is expected in a least-squares
approximation and does not imply that the computed minimizer is
inaccurate. The MSE quantifies the discrepancy between the data
and the fitted model, whereas the numerical indicators assess the
accuracy of the computational procedure.

The QR factorization produced a relative reconstruction error of
$7.3017\times10^{-17}$, while the orthogonality error was
$3.9468\times10^{-15}$. The least-squares optimality indicator,
$\|X^T(X\boldsymbol{\theta}^{*}-\mathbf{y})\|_2$, was
$2.3055\times10^{-14}$. These values are consistent with a
numerically accurate solution for this problem.

The uniqueness of the solution follows from the full column
rank of $X$. Indeed, if
$X\boldsymbol{\theta}=\mathbf{0}$, then the polynomial
$a x^2+bx+c$ vanishes at the six distinct values
$x_i=1,\ldots,6$. A nonzero polynomial of degree at most two
cannot have six distinct roots. Therefore,
$\boldsymbol{\theta}=\mathbf{0}$, so
$\operatorname{rank}(X)=3$ and the least-squares minimizer is
unique.

Figure~\ref{fig:quadratic-fit} presents the observations and the
fitted polynomial. The nearly straight fitted curve is consistent
with the numerically negligible quadratic coefficient.

\begin{figure}[htbp]
    \centering
    \includegraphics[width=0.82\textwidth]{local_results/problem2_quadratic_fit.pdf}
    \caption{Experimental observations and the quadratic least-squares fitted curve.}
    \label{fig:quadratic-fit}
\end{figure}

Since the observations are constructed from the enrollment digits
as prescribed by the assignment, the analysis concerns the
numerical behavior of the fitting procedure rather than an
empirical conclusion about predator--prey dynamics.